% ══════════════════════════════════════════════════════════════════════════════
\chapter{Correcteur orthographique}
\label{chap:td4}
% ══════════════════════════════════════════════════════════════════════════════

Avant d'interroger les index inversés, il faut s'assurer que les termes d'une requête
correspondent à des entrées connues du lexique. Un mot mal orthographié ne figurant
pas dans l'index retournerait zéro résultat, même si des documents pertinents existent.
Le correcteur orthographique développé en TD4 résout ce problème en appliquant un
pipeline déterministe à six étapes, combinant une heuristique de préfixe et la
distance de Levenshtein pour identifier la correction la plus vraisemblable.

\section{Distance de Levenshtein}

La \textbf{distance de Levenshtein} (cours, chapitre~3) mesure le nombre minimal
d'opérations élémentaires (insertion, suppression, substitution) pour transformer
une chaîne $a$ en une chaîne $b$. Elle est calculée par programmation dynamique :
\begin{equation}
  \operatorname{lev}(a, b) = \begin{cases}
    |a| & \text{si } |b| = 0 \\
    |b| & \text{si } |a| = 0 \\
    \operatorname{lev}(a_{1:}, b_{1:}) & \text{si } a_0 = b_0 \\
    1 + \min\!\left\{
      \operatorname{lev}(a_{1:}, b),\;
      \operatorname{lev}(a, b_{1:}),\;
      \operatorname{lev}(a_{1:}, b_{1:})
    \right\} & \text{sinon}
  \end{cases}
  \label{eq:levenshtein}
\end{equation}

L'implémentation optimise la mémoire en ne conservant que la ligne courante de la
matrice DP, réduisant l'empreinte de $O(|a|\cdot|b|)$ à $O(\min(|a|,|b|))$.

\begin{lstlisting}[caption={Implémentation Levenshtein optimisée en mémoire},
  label={lst:levenshtein}]
def levenshtein(a: str, b: str) -> int:
    if len(a) < len(b):
        return levenshtein(b, a)
    if not b:
        return len(a)
    prev = list(range(len(b) + 1))
    for i, ca in enumerate(a, 1):
        curr = [i]
        for j, cb in enumerate(b, 1):
            ins = prev[j] + 1
            dlt = curr[j - 1] + 1
            sub = prev[j - 1] + (ca != cb)
            curr.append(min(ins, dlt, sub))
        prev = curr
    return prev[-1]
\end{lstlisting}

\section{Pipeline de correction en six étapes}

La figure~\ref{fig:spell_pipeline} décrit le chemin de décision pour chaque terme
d'une requête.

\begin{figure}[ht]
  \centering
  \begin{tikzpicture}[
      box/.style={draw, rounded corners=3pt, fill=codegray,
                  minimum width=3.5cm, minimum height=0.65cm,
                  font=\small, align=center},
      dec/.style={draw, diamond, fill=white, aspect=2.5,
                  font=\small, align=center},
      arr/.style={-Stealth, thick},
      node distance=0.5cm and 1.3cm,
    ]
    \node[box] (start) {Terme de requête};
    \node[dec, below=of start] (ent) {Entité ?\\(nombre, date)};
    \node[box, right=1.4cm of ent] (ok_ent) {\texttt{entity}\\conservé tel quel};
    \node[dec, below=of ent] (exact) {Dans\\le lexique ?};
    \node[box, right=1.4cm of exact] (ok_exact) {\texttt{exact}\\retourne lemme};
    \node[dec, below=of exact] (cands) {$k$ candidats\\par préfixe};
    \node[box, right=1.4cm of cands] (nf) {\texttt{not\_found}\\introuvable};
    \node[box, below=0.9cm of cands] (lev) {Levenshtein\\$\to$ candidat min};
    \node[box, below=of lev] (out) {\texttt{CorrectionResult}\\(corrected, lemma)};

    \draw[arr] (start) -- (ent);
    \draw[arr] (ent) -- node[above, font=\tiny] {oui} (ok_ent);
    \draw[arr] (ent) -- node[left,  font=\tiny] {non} (exact);
    \draw[arr] (exact) -- node[above, font=\tiny] {oui} (ok_exact);
    \draw[arr] (exact) -- node[left,  font=\tiny] {non} (cands);
    \draw[arr] (cands) -- node[above, font=\tiny] {$k=0$} (nf);
    \draw[arr] (cands) -- node[left,  font=\tiny] {$k\geq1$} (lev);
    \draw[arr] (lev) -- (out);
  \end{tikzpicture}
  \caption{Pipeline de correction orthographique en six étapes : de la détection
    d'entité à la sélection du meilleur candidat par distance de Levenshtein.}
  \label{fig:spell_pipeline}
\end{figure}

La figure~\ref{fig:spell_pipeline} montre que le calcul de Levenshtein n'est déclenché
que lorsque la recherche par préfixe retourne plusieurs candidats ambigus, limitant le
coût de calcul aux cas réellement douteux. La recherche par préfixe utilise
\texttt{bisect.bisect\_left} sur la liste triée du lexique, offrant une complexité
$O(\log N + k)$ où $N$ est la taille du lexique et $k$ le nombre de candidats.

\section{Résultats et vulnérabilité identifiée}

\begin{table}[ht]
\centering
\caption{Exemples de corrections sur le corpus ADIT}
\label{tab:spell_examples}
\begin{tabular}{lllrc}
\toprule
\textbf{Terme saisi} & \textbf{Correction} & \textbf{Lemme retourné} & \textbf{dist.} & \textbf{Statut} \\
\midrule
\texttt{nanotechnolgie} & \texttt{nanotechnologie} & nanotechnologie & 1 & single\_candidate \\
\texttt{robotiqe}       & \texttt{robotique}       & robot           & 1 & single\_candidate \\
\texttt{algorithmf}     & \texttt{algorithme}      & algorithme      & 1 & best\_candidate   \\
\texttt{banotechnologie} & —                       & —               & — & not\_found        \\
\bottomrule
\end{tabular}
\end{table}

La \textbf{vulnérabilité principale} de l'heuristique est son aveuglement aux erreurs
portant sur le \emph{début} du mot. Comme l'illustre le tableau~\ref{tab:spell_examples},
\texttt{banotechnologie} (distance~1 de \texttt{nanotechnologie}) n'est pas retrouvé car
aucune entrée du lexique ne commence par \texttt{ban}. L'alternative robuste est
l'indexation par \textbf{trigrammes de caractères} : ces deux mots partagent 12 trigrammes
sur 13 (indice de Jaccard $\approx 0{,}92$), ce qui les rendrait voisins dans un index
trigram. Cette approche, utilisée en pratique par les moteurs de recherche industriels,
serait naturellement envisageable comme amélioration future. La couverture de tests
atteint 95,77\,\% sur 284 tests.
